\documentclass[a4paper]{article}

\usepackage{amsmath}
\usepackage{enumitem}

\newcommand{\entropy}[1]{\mathcal{H}(#1)}
\newcommand{\fnin}[2]{\((#1, #2)\)}

\begin{document}

\section{Addition}

For unsigned integer addition (i.e. with wrap-around semantics) of \(n\) bit
values, we prove that this operation retains half entropy by observing that
\begin{enumerate}[label=\alph*)]
    \item The output space of the addition is of size \(2^n\)
    \item The input space of the addition is of size \(2^{n+1}\)
    \item Each output is produced by \(2^{n}\) inputs
\end{enumerate}

The proof for (a) is trivial: there are a total of \(2^n\) representable
values, meaning the output space is at most \(2^n\).

The proof for (b) is also trivial: the input space of the addition is \(2^n *
2^n = 2^{n+1}\).

In order to prove (b), we start with the fact that output \(o\) can be obtained
by input \fnin{o}{0}. Now, we decrement the first value by 1, and increment the
second by 1, to obtain input \fnin{o - 1}{1}. The corresponding output is still
\(o\). We can repeat this operation up to \(2^n - 1\) times, as repeating it
for \(2^n\) times would yield \fnin{o - 2^n}{0 + 2^n} = \fnin{o}{0}, by
wrap-around semantics. We can repeat this logic for each original output \(o,
n\), where \(n \in [0, 2^n)\). Thus, each output value is mapped by \(2^n\)
input values.

Now, we simply apply the formula for Theil's U:

\begin{equation}
    U(I|O) = \frac{\entropy{I} - \entropy{I|O}}{\entropy{I}} = 1 - \frac{\entropy{I|O}}{\entropy{I}}
\end{equation}

For the input entropy, each input has \(p_i = \frac{1}{2^{n+1}}\). Thus:

\begin{equation*}
    \entropy{I} = \sum_{i} p_i * \log_2(p_i) = - 2^{n+1} * \frac{1}{2^{n+1}} * \log_2 (\frac{1}{2^{n+1}}) = n+1
\end{equation*}

For the conditional entropy, we have:

\begin{equation*}
    \entropy{I|O} = - \sum p(o, i) * \log_2 (\frac{p(o, i)}{p(o)}) = - 2^n * \frac{1}{2^n} * \log_2 ( \frac{\frac{1}{2^n}}{\frac{1}{2^n}} )
\end{equation*}

\end{document}
